\subsection{交错多线性映射丛}

在 7.8.1 至 7.8.5 号中，我们假设 K 的特征为 0，或所考虑的向量丛具有有限秩。尚不知道 7.8.1 中给出的函子 \(\alpha_{d}\) 的定义是否可以不加任何限制地进行。

7.8.1. 取 \(\mathbf{I}_{+} = \{0\}\)、\(\mathbf{I}_{-} = \{1\}\)，并令 \(d\) 为满足 \(\geqslant 1\) 的整数。我们定义一个

\(^{1}\) 当 \(M\) 为有限秩时，我们用 \(M^{*}\) 代替 \(M'\)。

向量函子 \(\alpha_{d}\)：将 \(\alpha_{d}(\mathcal{V})\) 记为由连续交错 \(d\)-线性映射构成的巴拿赫空间，这些映射从 \(\mathbf{V}_{1}^{d}\) 到 \(\mathbf{V}_{0}\)，并规定 \(\alpha_{d}(\mathbf{f})(u) = f_{0} \circ u \circ f_{1}^{d}\)，其中 \(u \in \alpha_{d}(\mathcal{V})\)。将向量丛 \(\alpha_{d}((\mathbf{M}_{1}, \mathbf{M}_{0}))\) 记为 \(\operatorname{Alt}^{d}(\mathbf{M}_{1}; \mathbf{M}_{0})\)，称为交错 \(d\)-线性映射向量丛，从 \(\mathbf{M}_{1}\) 到 \(\mathbf{M}_{0}\)。

从 \(\operatorname{Alt}^{d}(\mathbf{M}_{1};\mathbf{M}_{0})\) 到 \(\mathcal{L}(\mathbf{M}_{1},\ldots,\mathbf{M}_{1};\mathbf{M}_{0})\) 的典范嵌入是向量丛态射；\(\operatorname{Alt}^{d}(\mathbf{M}_{1};\mathbf{M}_{0})\) 是 \(\mathcal{L}(\mathbf{M}_{1},\ldots,\mathbf{M}_{1};\mathbf{M}_{0})\) 的子向量丛。

有 \(\operatorname{Alt}^1 (\mathbf{M}_1;\mathbf{M}_0) = \mathcal{L}(\mathbf{M}_1;\mathbf{M}_0)\)。置 \(\operatorname{Alt}^0 (\mathbf{M}_1;\mathbf{M}_0) = \mathbf{M}_0\)。

若 \(\omega\) 是一个截面 \(^{1}\)，属于 \(\mathsf{Alt}^{d}(\mathbf{M}_{1};\mathbf{M}_{0})\)，且 \(s_{1},\ldots,s_{d}\) 是 \(\mathbf{M}_{1}\) 的截面，则存在唯一一个 \(\mathbf{M}_{0}\) 的截面，记为 \(\omega(s_{1},\ldots,s_{d})\)，使得

\[
\omega (s _ {1}, \dots , s _ {d}) (b) = \omega (b) (s _ {1} (b), \dots , s _ {d} (b)) \quad \text { for all } b \in \mathbf {B}.
\]

7.8.2. 设 \(\varphi\) 是从 \(N \times_{B} N'\) 到 \(N''\) 的配对（参见第 7.3.1 号）；对每个 \(b\)，有一个双线性映射 \(\varphi_b\)，从 \(N_b \times N_b'\) 到 \(N_b''\)，它定义了（参见附录 A，第 III 节，第 142 页）一个从

\[
\operatorname{Alt} ^ {d} (\mathbf {M}; \mathbf {N}) _ {b} \times \operatorname{Alt} ^ {e} (\mathbf {M}; \mathbf {N} ^ {\prime}) _ {b}
\]

到 \(\operatorname{Alt}^{d + e}(\mathbf{M};\mathbf{N}^{\prime \prime})_{b}\) 的双线性映射。这些双线性映射组成一个配对 \(u\)，从

\[
\operatorname{Alt} ^ {d} (\mathbf {M}; \mathbf {N}) \times_ {\mathbf {B}} \operatorname{Alt} ^ {e} (\mathbf {M}; \mathbf {N} ^ {\prime})
\]

到 \(\mathsf{Alt}^{d + e}(\mathbf{M};\mathbf{N}^{\prime \prime})\)；若 \(\omega\) 和 \(\omega^{\prime}\) 分别是 \(\mathsf{Alt}^{d}(\mathbf{M};\mathbf{N})\) 和 \(\mathsf{Alt}^{e}(\mathbf{M};\mathbf{N}^{\prime})\) 在开集 U 上的截面，则 \(u(\omega ,\omega^{\prime})\) 是 \(\mathsf{Alt}^{d + e}(\mathbf{M};\mathbf{N}^{\prime \prime})\) 在 U 上的截面，记为 \(\omega \wedge_{\varphi}\omega^{\prime}\)，称为 \(\omega\) 与 \(\omega^{\prime}\) 的外积。

有公式：

\[
(1) \quad \left(\omega \wedge_ {\varphi} \omega^ {\prime}\right) \left(s _ {1}, \dots , s _ {d + e}\right) = \sum_ {\sigma} \varepsilon_ {\sigma} \varphi \left(\omega \left(s _ {\sigma (1)}, \dots , s _ {\sigma (d)}\right), \omega^ {\prime} \left(s _ {\sigma (d + 1)}, \dots , s _ {\sigma (d + e)}\right)\right)
\]

其中 \(s_{i}\) 是 M 在 U 上的截面，并且求和遍历置换 \(\sigma\)，这些置换属于 \(\{1,2,\ldots,d+e\}\)，且满足以下条件：

\[
\sigma (1) <   \dots <   \sigma (d) \quad \text { and } \quad \sigma (d + 1) <   \dots <   \sigma (d + e).
\]

7.8.3. 设 M 是一个向量丛，A 是一个以 B 为基底的代数丛。假设 A 的纤维 \(A_{b}\) 是结合且交换的代数，并带有单位元，记为 \(e_{b}\)。对 B 的每个开集 U，记 \(\Omega^{d}(U)\) 为其截面所成的 \(\mathcal{C}^{r}(U)\)-模，其中这些截面属于丛 \(\operatorname{Alt}^{d}(M;A)\)，并记 \(\Omega^{*}(U)\) 为 \(\Omega^{d}(U)\) 的直和，其中 \(d\geqslant0\)。每个纤维上的乘法定义了从 \(A\times_{B}A\) 到 A 的配对，因而由 7.8.2 在 \(\Omega^{*}(U)\) 上得到一个分次代数结构，它是结合且反交换的。子代数 \(\Omega^{0}(U)\) 是 A 的截面代数。一个元素 \(\omega\) 属于 \(\Omega^{1}(\mathbf{U})\) 时，可将其视为从 \(\mathbf{M}|\mathbf{U}\) 到 \(\mathbf{A}|\mathbf{U}\) 的 U-态射（7.7.3）：若 \(s\in\mathcal{S}_{\mathbf{M}}^{r}(\mathbf{U})\)，则将 \(\langle\omega,s\rangle\) 记为 A 的截面 \(\omega(s)\)（7.4.2）。设 \(s_{j}\in\mathcal{S}_{\mathbf{M}}^{r}(\mathbf{U})\) 且 \(\omega_{j}\in\Omega^{1}(\mathbf{U})\)（\(1\leqslant j\leqslant d\)）；有：

\(^{1}\) 读者应注意不要把这里对字母 \(\omega\) 的用法与第 10 页定义的用法混淆。

\[
\omega (s _ {1}, \dots , s _ {d}) = \det (\langle \omega_ {i}, s _ {j} \rangle) \quad \text {   for   } \omega = \omega_ {1} \wedge \dots \wedge \omega_ {d}.\tag{2}
\]

7.8.4. 设 \(d \geqslant 1\)。存在一个配对 \(i\)，从 \(\mathbf{M} \times_{\mathbf{B}} \mathrm{Alt}^{d}(\mathbf{M}; \mathbf{A})\) 到 \(\mathrm{Alt}^{d-1}(\mathbf{M}; \mathbf{A})\)，其在每个纤维上的限制由右内积给出（参见附录 A，第 III 节，第 156 页）。若 \(s\) 是 \(\mathbf{M}\) 在开子集 \(\mathbf{U}\) 上的截面，且 \(\omega \in \Omega^{d}(\mathbf{U})\)，则将 \(i(s)\omega\) 记为截面 \(i(s, \omega)\)，它属于 \(\mathrm{Alt}^{d-1}(\mathbf{M}; \mathbf{A})\) 且定义在 \(\mathbf{U}\) 上；置 \(i(s)\omega = 0\)，其中 \(\omega\) 属于 \(\Omega^{0}(\mathbf{U})\)。

因此，对每个截面 \(s\)，它是 \(\mathbf{M}\) 在 \(\mathbf{U}\) 上的截面，我们关联一个 \(\mathcal{C}^r (\mathbf{U})\)-模 \(\Omega^{*}(\mathbf{U})\) 的自同态。公式如下：

\[
(3) \quad (i (s) \omega) (s _ {1}, \dots , s _ {d - 1}) = \omega (s, s _ {1}, \dots , s _ {d - 1}) \quad \text { for } \omega \in \Omega^ {d} (\mathbf {U}), d \geqslant 1
\]

\begin{enumerate}
\def\labelenumi{(\arabic{enumi})}
\setcounter{enumi}{3}
\tightlist
\item
  \(i(s)\circ i(s) = 0\)
\end{enumerate}

\[
i (s) \omega = \langle \omega , s \rangle \quad \text { for } \omega \in \Omega^ {1} (\mathbf {U}) \tag {5}
\]

\[
(6) \quad i (s). (\omega \wedge \omega^ {\prime}) = i (s) \omega \wedge \omega^ {\prime} + (- 1) ^ {d} \omega \wedge i (s) \omega^ {\prime} \quad \text { for } \omega \in \Omega^ {d} (\mathbf {U})
\]

\[
i (s) (\omega_ {1} \wedge \dots \wedge \omega_ {p}) = \sum_ {i = 1} ^ {p} (- 1) ^ {i + 1} \langle \omega_ {i}, s \rangle \omega_ {1} \wedge \dots \wedge \hat {\omega} _ {i} \wedge \dots \wedge \omega_ {d}.\tag{7}
\]

在最后一个公式中，\(\omega_{i}\) 属于 \(\Omega^{1}(\mathbf{U})\)，上置的符号表示应将其所标出的符号略去。

上面对截面所描述的所有运算都是关于环 \(\mathcal{C}^r (\mathbf{U};\mathbf{K})\) 的多线性的。

7.8.5. 设 L 是 K 上的巴拿赫代数。7.7 和 7.8 节的定义和结果可推广到 L 上的向量丛情形：类似地定义 L-多线性映射丛或交错 L-多线性映射丛。

